\section{Supplementary Derivations for the Token Cost Model}
\label{sec:cost-derivations}

In this section, we provide detailed derivations and proofs for formulas and results used directly in Section~\ref{sec:cost}. As in the main text, the number of listwise calls is used as a linear proxy of total token cost. The total merge-stage token cost is the sum of the per-pass costs in Eq.~\eqref{eq:total-merge-cost}; the following proofs establish the common number of merge passes and the per-pass terms for each algorithm.

\subsection{Number of Merge Passes}

\begin{proposition}[$k$-way Merge Token Cost]
For $n>m$ and $k\ge2$, all three merge algorithms require
\begin{align}
  \Pass_a(k)
  &=\left\lceil
    \log_k\left\lceil\frac{n}{m}\right\rceil
    \right\rceil,
    \qquad a\in\{\SD,\LT,\SLT\}.
\end{align}
\end{proposition}

\begin{proof}
Let $r_t$ be the number of runs at the beginning of merge pass $t$. Phase~1 produces $r_0=\lceil n/m\rceil$ runs. Since one $k$-way merge pass produces one output run for each group of at most $k$ input runs, $r_{t+1}=\lceil r_t/k\rceil$. For every positive $x$ and integer $k\ge1$, the nested ceiling identity $\lceil\lceil x\rceil/k\rceil=\lceil x/k\rceil$ holds. Applying this identity to the recurrence and inducting on $t$ yields $r_t=\lceil r_0/k^t\rceil$. The merge terminates at the least $t$ for which $r_t\le1$, equivalently $k^t\ge r_0$. This least integer is $\lceil\log_k r_0\rceil$, which proves Eq.~\eqref{eq:number_merge_pass}.
\end{proof}

\subsection{Per-pass Token Cost for Sliding Window}

For a listwise call in which all $k$ runs are active, each run contributes $h_k=\lfloor m/k\rfloor$ items, giving an actual batch size of $b_k=kh_k\le m$; the feasible range $2\le k\le\lfloor m/2\rfloor$ ensures $h_k\ge2$. We assume uniform random interleaving so that no run is systematically exhausted early; otherwise, premature exhaustion of a particular run's $h_k$ items would shorten the safe prefix and reduce throughput.

\begin{theorem}[$\Pcost_{\SD}$ for sliding window]
Under uniform random interleaving, the expected safe prefix length of a listwise call is
\begin{align}
  \mu_{\SD}(k,m)
  &=(b_k+1)\prod_{\ell=1}^{k}
    \frac{\ell h_k}{\ell h_k+1}.
\end{align}
\end{theorem}

\begin{proof}
Assign $h_k$ independent uniform points to each run, and let $M_i$ be the largest point assigned to run $i$. For $x\in[0,1]$, the distribution of a run maximum and the survival function of the first exhaustion threshold $\tau=\min_i M_i$ are
\begin{align}
  &\Pr(M_i\le x) = x^{h_k}\\
  \Pr(\tau>x) &= \prod_{i=1}^{k}\Pr(M_i>x)=(1-x^{h_k})^k
\end{align}
Let $Y_k$ be the emitted safe prefix length. The run labels are independent of the ordered point locations, so conditional on $Y_k=j$, the threshold $\tau$ is the $j$th order statistic among $b_k$ uniform points and has mean $j/(b_k+1)$. Therefore,
\begin{align}
  &\mathbb E[\tau\mid Y_k] =\frac{Y_k}{b_k+1}\\
  &\mathbb E[Y_k] =(b_k+1)\mathbb E[\tau]
\end{align}
Using the tail integral identity, substituting $u=x^{h_k}$, and applying $\Gamma(z+1)=z\Gamma(z)$ gives
\begin{align}
  \mathbb E[Y_k]
  &=(b_k+1)\int_0^1(1-x^{h_k})^kdx
  \nonumber\\
  &=\frac{b_k+1}{h_k}
    \int_0^1u^{1/h_k-1}(1-u)^kdu
    \nonumber\\
  &=\frac{b_k+1}{h_k}
    B\!\left(\frac1{h_k},k+1\right)
    \nonumber\\
  &=(b_k+1)
    \frac{\Gamma(1+1/h_k)\Gamma(k+1)}
         {\Gamma(k+1+1/h_k)}
    \nonumber\\
  &=(b_k+1)\prod_{\ell=1}^{k}
    \frac{\ell h_k}{\ell h_k+1}.
\end{align}
This is Eq.~\eqref{eq:mu-sw}.
\end{proof}

\subsection{Per-round Token Cost for Loser Tree}

\begin{proposition}[$\Pcost_{\LT}$ for loser tree]
The per-round token cost at the t-th merge pass for loser tree merge is
\begin{align}
  \Pcost_{\LT}(t,k;n,m)
  &=n+G_t\left(
      \Cost_{\SD}(2;k,m)
      +\left\lceil\frac{k}{m}\right\rceil
    \right).
\end{align}
\end{proposition}

\begin{proof}
For each of the $G_t$ merge groups, tree initialization ranks the $k$ representatives in $\lceil k/m\rceil$ Phase~1 listwise calls and then uses $\Cost_{\SD}(2;k,m)$ calls for a two-way sliding window merge. Since initialization typically produces two chunks in our setting, this merge uses fan-in two and is not optimized separately; the same setting is used for segmented loser tree initialization. During the merge, each emitted item triggers one \textsc{RefreshPath} call, giving $n$ path refresh calls over all groups. Adding these $n$ refresh calls and the $G_t$ identical initialization costs proves Eq.~\eqref{eq:pcost-lt}.
\end{proof}

Here, $n$ denotes the total number of items and typically exceeds 100, whereas the initialization term $G_t\bigl(\Cost_{\SD}(2;k,m)+\lceil k/m\rceil\bigr)$ contributes only a small number of listwise calls. Hence, $n\gg G_t\bigl(\Cost_{\SD}(2;k,m)+\lceil k/m\rceil\bigr)$, and the selector in Section~\ref{sec:cost} uses the approximation $\Pcost_{\LT}(t,k;n,m)=n$. The total cost is therefore approximately $n\Pass_{\LT}(k)$, which is linear in the number of merge passes with slope $n$. It is minimized when $\Pass_{\LT}(k)=1$, so all runs are merged in a single pass using one loser tree.

\subsection{Per-round Token Cost for Segmented Loser Tree}

\begin{proposition}[$\Pcost_{\SLT}$ for segmented loser tree]
The per-round token cost at the t-th merge pass for segmented loser tree merge is
\begin{align}
  \Pcost_{\SLT}(t,k;n,m)
  &=\left\lceil\frac{n}{h}\right\rceil
    +G_t\left(
      \Cost_{\SD}(2;kh,m)
      +\left\lceil\frac{kh}{m}\right\rceil
    \right).
\end{align}
\end{proposition}

\begin{proof}
For each of the $G_t$ merge groups, tree initialization ranks the $kh$ representatives in $\lceil kh/m\rceil$ Phase~1 listwise calls and then uses $\Cost_{\SD}(2;kh,m)$ calls for the same two-way sliding window merge. Each path refresh processes $h$ items. In the cost model, the item counts of the merge groups are summed before the ceiling is applied, giving $\lceil n/h\rceil$ path refresh calls over the pass. Adding this term to the $G_t$ identical initialization costs proves Eq.~\eqref{eq:pcost-slt}.
\end{proof}

\section{Algorithm Details}
\label{sec:appendix-algorithms}

\subsection{Loser Tree Construction and Path Refresh}

Each internal node stores the loser of its most recent match; the root additionally exposs the current winner.  

\begin{algorithm}[t]
\small
\caption{\textsc{InitLoserTree}}
\label{alg:init_loser_tree}
\KwIn{active run indices $A$, representatives $\{x_i\mid i\in A\}$,
      instruction $\phi$, batch capacity $m$}
\KwOut{a binary loser tree $T$}

$Q_{\mathrm{init}}\gets
  \textsc{FullSlidingWindowMerge}(\{[x_i]\mid i\in A\},\phi,m)$\;
Construct $\rho$ such that $\rho(x)$ is the position of $x$ in
$Q_{\mathrm{init}}$\;
Construct an exact binary tree $T$ with one leaf for every $i\in A$\;
\ForEach{$i\in A$}{$\textsc{SetWinner}(\textsc{Leaf}(i),i)$\;}
\ForEach{internal node $v$ in bottom-up order}{
  $a\gets\textsc{CurrentWinner}(\textsc{LeftChild}(v))$\;
  $b\gets\textsc{CurrentWinner}(\textsc{RightChild}(v))$\;
  \uIf{$\rho(x_a)<\rho(x_b)$}{
    $\textsc{SetWinner}(v,a)$;$\textsc{SetStoredLoser}(v,b)$\;
  }
  \Else{
    $\textsc{SetWinner}(v,b)$;$\textsc{SetStoredLoser}(v,a)$\;
  }
}
$\textsc{SetCurrentWinner}(T,\textsc{CurrentWinner}(\textsc{Root}(T)))$\;
\Return{$T$}\;
\end{algorithm}

\begin{algorithm}[t]
\small
\caption{\textsc{RefreshPath}}
\label{alg:refresh_loser_tree_path}
\KwIn{loser tree $T$, changed run index $u$, active indices $A$,
      representatives $\{x_i\mid i\in A\}$, instruction $\phi$}
\KwOut{updated loser tree $T$}

$B\gets[\ ]$\;
\lIf{$u\in A$}{Append $x_u$ to $B$}
$z\gets\textsc{Leaf}(u)$\;
\While{$z\ne\textsc{Root}(T)$}{
  $p\gets\textsc{Parent}(z)$; $q\gets\textsc{StoredLoser}(p)$\;
  \lIf{$q\in A$ \textbf{and} $x_q\notin B$}{Append $x_q$ to $B$}
  $z\gets p$\;
}
\If{$B\ne\emptyset$}{
  $\mathcal O_{\mathrm{path}}\gets\textsc{ListwiseRank}(B,\phi)$\;
  Construct $\rho_{\mathrm{path}}$ from $\mathcal O_{\mathrm{path}}$\;
}
$z\gets\textsc{Leaf}(u)$\;
$a\gets u$ if $u\in A$, and $a\gets\textsc{null}$ otherwise\;
\While{$z\ne\textsc{Root}(T)$}{
  $p\gets\textsc{Parent}(z)$; $q\gets\textsc{StoredLoser}(p)$\;
  \lIf{$q\notin A$}{$q\gets\textsc{null}$}
  \uIf{$a=\textsc{null}$ \textbf{or} $q=\textsc{null}$}{
    $\textsc{SetStoredLoser}(p,\textsc{null})$;
    $a\gets q$ if $a=\textsc{null}$\;
  }
  \uElseIf{$\rho_{\mathrm{path}}(x_a)<\rho_{\mathrm{path}}(x_q)$}{
    $\textsc{SetStoredLoser}(p,q)$\;
  }
  \Else{
    $\textsc{SetStoredLoser}(p,a)$;$a\gets q$\;
  }
  $\textsc{SetWinner}(p,a)$;$z\gets p$\;
}
$\textsc{SetCurrentWinner}(T,a)$\;
\Return{$T$}\;
\end{algorithm}


\subsection{Segmented Loser Tree and Pointer-Guided Path Refresh}
\label{sec:appendix-segment-algorithm}

The segmented loser tree applies each loser-tree match to $h$-item
segments.  An internal node $v$ stores only
$\mathit{Loser}(v)$ and $\mathit{Pointer}(v)$; the winner $tmp$ is a
transient segment passed to its parent.  For a group of $k$ runs, we set
\begin{equation}
  h=\left\lfloor
    \frac{m}{1+\lceil\log_2 k\rceil}
  \right\rfloor,
  \label{eq:segment-length}
\end{equation}
so that one replacement segment and the losers on a root-to-leaf path
fit in a single listwise call.  $\textsc{FillLeaf}(v,h)$ returns the next
segment from the run associated with leaf $v$.

\begin{algorithm}[t]
\small
\caption{Segmented Loser Tree Merge}
\label{alg:segmented_loser_tree_merge}
\KwIn{a group of sorted runs $\mathcal G=\{R_1,\ldots,R_k\}$,
      instruction $\phi$}
\KwOut{the merged order $\Pi$}

$\Pi\gets[\ ]$\;

Construct a binary segmented loser tree $T$ over $\mathcal G$
with segment length $h$\;

\algocomment{load and jointly rank the initial leaf segments}\;
\ForEach{leaf $v$ concurrently}{
  $S_v\gets\textsc{FillLeaf}(v,h)$\nllabel{line:slt-init}\;
}
Rank all $S_v$ with \textsc{FullSlidingWindowMerge}\;
Let $\rho_0$ be the resulting rank map\;

\algocomment{initialize internal-node states using the shared rank map}\;
\Fn{\textsc{BuildSubtree}$(v,\rho_0)$}{
  \If{$v$ is a leaf}{\Return{$S_v$}\;}
  $L\gets\textsc{BuildSubtree}(\textsc{LeftChild}(v),\rho_0)$\;
  $R\gets\textsc{BuildSubtree}(\textsc{RightChild}(v),\rho_0)$\;
  \Return{$\textsc{ReplaySegmentNode}(v,L,R,\rho_0)$}
    \nllabel{line:slt-build}\;
}

\algocomment{emit winner segments using pointer-guided path refresh}\;
$tmp\gets\textsc{BuildSubtree}(\textsc{Root}(T),\rho_0)$\;
\While{$tmp\ne\emptyset$}{
  Append $tmp$ to $\Pi$\nllabel{line:slt-emit}\;
  $tmp\gets\textsc{PointerGuidedPathRefresh}(T,\phi)$
    \nllabel{line:slt-repair}\;
}
\Return{$\Pi$}\;
\end{algorithm}

% \begin{algorithm}[t]
% \small
% \linespread{0.88}\selectfont
% \caption{Segmented Loser Tree Merge}
% \label{alg:segmented_loser_tree_merge}
% \KwIn{a group of sorted runs $\mathcal G=\{R_1,\ldots,R_g\}$,
%       instruction $\phi$, listwise batch size $m$, optional output limit $K$}
% \KwOut{the first $K$ items of the merged order $\Pi$; the full order if
%        $K=\infty$}

% Remove empty runs from $\mathcal G$\;
% $g\gets|\mathcal G|$\;
% $\Pi\gets[\ ]$\;
% \If{$g=0$}{\Return{$\Pi$}\;}
% \If{$g=1$}{\Return{the first $K$ items of $R_1$}\;}
% $d\gets\max(1,\lceil\log_2 g\rceil)$\;
% $h\gets\max(1,\lfloor m/(1+d)\rfloor)$\;
% Require $(1+d)h\le m$\;
% Construct a binary tree $T$ with exactly one leaf for every
% run in $\mathcal G$\;
% Set $\mathit{done}(i)\gets\textsc{false}$ for every leaf run $i$\;
% \ForEach{node $v$ in $T$}{
%   $\mathit{ready}(v)\gets[\ ]$\;
%   $\mathit{residual}(v)\gets[\ ]$\;
%   $\mathit{blocked}(v)\gets\textsc{null}$\;
% }
% \ForEach{leaf $v$ concurrently}{
%   $\textsc{FillLeaf}(v,h)$\nllabel{line:slt-init}\;
% }
% $Q_0\gets\textsc{FullSlidingWindowMerge}
%   (\{\mathit{ready}(v)\mid v\text{ is a leaf}\},\phi,m)$\;
% Construct $\rho_0$ such that $\rho_0(x)$ is the position of $x$ in $Q_0$\;
% \ForEach{internal node $v$ in bottom-up order}{
%   $L\gets\textsc{TakeReady}(\textsc{LeftChild}(v),h)$\;
%   $R\gets\textsc{TakeReady}(\textsc{RightChild}(v),h)$\;
%   $\textsc{ReplaySegmentNode}(v,L,R,\rho_0,[\ ])$
%   \nllabel{line:slt-build}\;
% }
% $r\gets\textsc{Root}(T)$\;
% \While{$|\Pi|<K$ \textbf{and} \textsc{HasMore}$(r)$}{
%   \If{$\mathit{ready}(r)=\emptyset$}{
%     $ok\gets\textsc{PointerGuidedPathRefresh}(T,\phi,m,h)$
%     \nllabel{line:slt-repair}\;
%     \lIf{\textbf{not} $ok$}{\textbf{break}}
%   }
%   $E\gets\textsc{TakeReady}(r,\min(h,K-|\Pi|))$\;
%   Append $E$ to $\Pi$\nllabel{line:slt-emit}\;
% }
% \Return{$\Pi$}\;
% \end{algorithm}
The rank map $\rho$ defines the order of $L\uplus R$.
Membership $x\in L$ or $x\in R$ refers to the replay input supplying
$x$, not merely its value.  Algorithm~\ref{alg:replay_segment_node}
shows the common case $|L|=|R|=h$; final partial segments are handled
analogously using their actual lengths.
\begin{algorithm}[t]
\small
\caption{\textsc{ReplaySegmentNode}}
\label{alg:replay_segment_node}
\KwIn{node $v$, $h$-item replay segments $L,R$ associated with its
      left and right child slots, rank map $\rho$}
\KwOut{winner segment $tmp$; updated
       $\mathit{Loser}(v)$ and $\mathit{Pointer}(v)$}

$Z\gets\operatorname{Sort}_{\rho}(L\uplus R)$
  \nllabel{line:replay-merge}\;

$n_L,n_R\gets0,0$\;
\ForEach{$x\in Z$ in order}{
  \uIf{$x\in L$}{
    $n_L\gets n_L+1$\;
  }
  \Else{
    $n_R\gets n_R+1$\;
  }
  \lIf{$n_L=h$ \textbf{or} $n_R=h$}{\textbf{break}}
}
$\mathit{Pointer}(v)\gets
  \begin{cases}
    \textsc{LeftChild}(v),  & n_L=h,\\
    \textsc{RightChild}(v), & n_R=h
  \end{cases}$
  \nllabel{line:replay-split}\;

$tmp\gets Z[1..h]$
  \nllabel{line:replay-winner}\;
$\mathit{Loser}(v)\gets Z[h+1..2h]$
  \nllabel{line:replay-loser}\;

\Return{$tmp$}\;
\end{algorithm}

% \begin{algorithm}[t]
% \small
% \linespread{0.88}\selectfont
% \caption{\textsc{ReplaySegmentNode}}
% \label{alg:install_segment_state}
% \KwIn{internal node $v$, certified left and right child segments $L,R$,
%       rank map $\rho$, certified prefix $P$ already stored at $v$}
% \KwOut{$v$ with an updated certified prefix, retained segment, and
%        blocked-child pointer}

% $\ell\gets\textsc{LeftChild}(v)$\;
% $r\gets\textsc{RightChild}(v)$\;
% $more_\ell\gets\textsc{HasMore}(\ell)$\;
% $more_r\gets\textsc{HasMore}(r)$\;
% $\mathit{ready}(v)\gets P$\;
% $\mathit{residual}(v)\gets[\ ]$\;
% $\mathit{blocked}(v)\gets\textsc{null}$\;
% \uIf{$L=\emptyset$ \textbf{and} $R=\emptyset$}{
%   Require $\neg(more_\ell\wedge more_r)$
%   \nllabel{line:install-empty-start}\;
%   \uIf{$more_\ell$}{$\mathit{blocked}(v)\gets\ell$}
%   \ElseIf{$more_r$}{$\mathit{blocked}(v)\gets r$}
% }
% \uElseIf{$L=\emptyset$}{
%   \uIf{$more_\ell$}{
%     $\mathit{residual}(v)\gets R$\;
%     $\mathit{blocked}(v)\gets\ell$\;
%   }
%   \Else{
%     Append $R$ to $\mathit{ready}(v)$\;
%     \lIf{$more_r$}{$\mathit{blocked}(v)\gets r$}
%   }
% }
% \uElseIf{$R=\emptyset$}{
%   \uIf{$more_r$}{
%     $\mathit{residual}(v)\gets L$\;
%     $\mathit{blocked}(v)\gets r$\;
%   }
%   \Else{
%     Append $L$ to $\mathit{ready}(v)$\;
%     \lIf{$more_\ell$}{$\mathit{blocked}(v)\gets\ell$\nllabel{line:install-empty-end}}
%   }
% }
% \Else{
%   $Z\gets\textsc{StableMerge}_{\rho}(L,R)$
%   \nllabel{line:install-merge}\;
%   \algocomment{preserve the order inside $L$ and $R$} \\
%   Let $b$ be the first position in $Z$ at which $L$ or $R$ is exhausted\;
%   Let $e$ be that exhausted child and $o$ the other child\;
%   \uIf{$\textsc{HasMore}(e)$}{
%     Append $Z[1..b]$ to $\mathit{ready}(v)$\;
%     $\mathit{residual}(v)\gets Z[b+1..|Z|]$\;
%     $\mathit{blocked}(v)\gets e$\nllabel{line:install-barrier}\;
%   }
%   \Else{
%     Append $Z$ to $\mathit{ready}(v)$\;
%     \lIf{$\textsc{HasMore}(o)$}{$\mathit{blocked}(v)\gets o$\nllabel{line:install-exhausted}}
%   }
% }
% \Return{the updated state of $v$}\;
% \end{algorithm}
\noindent\textbf{Initialization.}
Algorithm~\ref{alg:segmented_loser_tree_merge} jointly ranks the initial
leaf segments.  \textsc{BuildSubtree} recursively passes each winner
upward, while every internal node retains only its loser and pointer.
The winner returned by the root becomes the first output segment.

\noindent\textbf{Pointer-guided refresh.}
After a root winner is emitted, Algorithm~\ref{alg:repair_segment_path}
follows the pointers to the affected leaf and jointly ranks
\[
  \mathcal B
  =
  tmp\uplus\biguplus_{v\in\mathcal P}\mathit{Loser}(v).
\]
During bottom-up replay, the old pointer places $tmp$ and the stored loser
in their replay slots.  \textsc{ReplaySegmentNode} then updates the node
state and passes the next transient winner toward the root.
\begin{algorithm}[t]
\small
\caption{Pointer-Guided Path Refresh}
\label{alg:repair_segment_path}
\KwIn{segmented loser tree $T$, instruction $\phi$}
\KwOut{the next $h$-item winner segment $tmp$ at the root}

$\mathcal P\gets[\ ]$; $u\gets\textsc{Root}(T)$\;

\algocomment{follow the pointers to the affected leaf}\;
\While{$u$ is internal}{
  Append $u$ to $\mathcal P$\;
  $u\gets\mathit{Pointer}(u)$
    \nllabel{line:repair-descend}\;
}
$tmp\gets\textsc{FillLeaf}(u,h)$\;

\algocomment{jointly rank the segments on the refresh path}\;
$\mathcal B\gets
  tmp\uplus\biguplus_{v\in\mathcal P}\mathit{Loser}(v)$\;
$\mathcal O_{\mathcal B}
  \gets\textsc{ListwiseRank}(\mathcal B,\phi)$
  \nllabel{line:repair-rank}\;
Let $\rho_{\mathcal B}$ be the rank map induced by
$\mathcal O_{\mathcal B}$\;

\algocomment{replay the winner segment toward the root}\;
\ForEach{$v\in\mathcal P$ in bottom-up order}{
  \uIf{$\mathit{Pointer}(v)=\textsc{LeftChild}(v)$}{
    $L\gets tmp$; $R\gets\mathit{Loser}(v)$\;
  }
  \Else{
    $L\gets\mathit{Loser}(v)$; $R\gets tmp$\;
  }
  $tmp\gets
    \textsc{ReplaySegmentNode}(v,L,R,\rho_{\mathcal B})$
    \nllabel{line:repair-replay}\;
}

\Return{$tmp$}\;
\end{algorithm}

% \begin{algorithm}[t]
% \small
% \linespread{0.88}\selectfont
% \caption{Pointer-Guided Path Refresh}
% \label{alg:repair_segment_path}
% \KwIn{segmented loser tree $T$, instruction $\phi$, listwise batch size $m$,
%        segment length $h$}
% \KwOut{updated $T$; \textsc{true} iff the root has a nonempty certified prefix}

% $r\gets\textsc{Root}(T)$\;
% \If{$\mathit{ready}(r)\ne\emptyset$}{\Return{\textsc{true}}\;}
% $P\gets[\ ]$\;
% $u\gets r$\;
% \While{$u$ is not a leaf}{
%   $b\gets\mathit{blocked}(u)$\;
%   \If{$b=\textsc{null}$}{\Return{\textsc{false}}\;}
%   Append $u$ to $P$\;
%   $u\gets b$\nllabel{line:repair-descend}\;
%   \If{$u$ is internal \textbf{and} $\mathit{blocked}(u)=\textsc{null}$}{
%     \textbf{break}\;
%   }
% }
% \ForEach{$v\in P$}{
%   Save $Q_v\gets\mathit{ready}(v)$,
%        $Z_v\gets\mathit{residual}(v)$, and
%        $b_v\gets\mathit{blocked}(v)$\nllabel{line:repair-save}\;
% }
% $c\gets\sum_{v\in P}(|Q_v|+|Z_v|)$\;
% Require $c+h\le m$\;
% \If{$u$ is a leaf}{\textsc{FillLeaf}$(u,h)$\;}
% $I\gets\textsc{TakeReady}(u,h)$\;
% $B\gets I$ followed by all $Q_v$ and $Z_v$ for $v\in P$\;
% $\mathcal O_B\gets\textsc{ListwiseRank}(B,\phi)$
% \nllabel{line:repair-rank}\;
% Construct $\rho_B$ such that $\rho_B(x)$ is the position of $x$ in
% $\mathcal O_B$\;
% $through\gets I$\;
% $child\gets u$\;
% \ForEach{$v\in P$ in bottom-up order}{
%   Require $b_v=child$\;
%   \uIf{$child=\textsc{LeftChild}(v)$}{
%     $L\gets through$; $R\gets Z_v$\;
%   }
%   \Else{
%     $L\gets Z_v$; $R\gets through$\;
%   }
%   $\textsc{ReplaySegmentNode}(v,L,R,\rho_B,Q_v)$
%   \nllabel{line:repair-replay}\;
%   $child\gets v$\;
%   \If{$v\ne r$}{
%     $through\gets\textsc{TakeReady}(v,h)$
%     \nllabel{line:repair-propagate}\;
%   }
% }
% \Return{$\mathit{ready}(r)\ne\emptyset$}\;
% \end{algorithm}


\subsection{Pipelined Execution for \texorpdfstring{\textsf{LIMIT}~$K$}{LIMIT K}}
\label{sec:topk-pseudocode}

Every internal merge producer is driven by a bounded background pump.  The default look-ahead queue has capacity one.  Closing the root cancels its pump and recursively closes all upstream producers, so early termination bounds speculative work.

\begin{algorithm}[t]
\small
\caption{Pipelined Top-$K$ Merge}
\label{alg:pipelined_topk_driver}
\KwIn{Phase-1 runs $\mathcal L=[R_1,\ldots,R_r]$, instruction $\phi$,
      batch size $m$, fan-in $k$, limit $K$, group-merger type}
\KwOut{semantic top-$K$ prefix $\Pi_K$}

$S\gets[\textsc{QueueProducer}(R_1),\ldots,
         \textsc{QueueProducer}(R_r)]$\;
\While{$|S|>1$}{
  Partition $S$ into groups $\mathcal G_1,\ldots,\mathcal G_q$ of at most
  $k$ producers\;
  $S_{\mathrm{next}}\gets[\ ]$\;
  \ForEach{$\mathcal G_j$}{
    \eIf{$|\mathcal G_j|=1$}{
      Append its only producer to $S_{\mathrm{next}}$\;
    }{
      Create a pumped sliding-window, binary loser-tree, or segmented
      loser-tree producer $M_j$ over $\mathcal G_j$ with budget $K$ and
      append it to $S_{\mathrm{next}}$\;
    }
  }
  $S\gets S_{\mathrm{next}}$\;
}
$M_{\mathrm{root}}\gets S[1]$; $\Pi_K\gets[\ ]$\;
\While{$|\Pi_K|<K$ \textbf{and} $M_{\mathrm{root}}$ is not exhausted}{
  $X\gets\textsc{Next}(M_{\mathrm{root}},K-|\Pi_K|)$\;
  \lIf{$X=\emptyset$}{\textbf{break}}
  Append $X$ to $\Pi_K$\;
}
$\textsc{Close}(M_{\mathrm{root}})$; remove temporary ids from $\Pi_K$\;
\Return{$\Pi_K$}\;
\end{algorithm}

\begin{algorithm}[t]
\small
\caption{$\textsc{Next}(M,c)$ for a Sliding-Window Producer}
\label{alg:pipelined_sw_next}
\KwIn{producer $M$ over children $S_1,\ldots,S_g$, request $c$,
      batch size $m$, instruction $\phi$, budget $K$}
\KwOut{up to $c$ next items}

$c\gets\min(c,K-\mathit{delivered}(M))$\;
\While{$|\mathit{ready}(M)|<c$ \textbf{and} $M$ is not exhausted}{
  $A\gets$ children with buffered items or remaining upstream supply\;
  \lIf{$A=\emptyset$}{\textbf{break}}
  \If{$|A|=1$}{
    Drain that child's remaining stream into $\mathit{ready}(M)$, up to the
    remaining budget, then \textbf{break}\;
  }
  $h\gets\lfloor m/|A|\rfloor$\;
  Pull from all children in $A$ concurrently until each contributes up to
  $h$ items to the sticky window\;
  $\mathcal O\gets\textsc{ListwiseRank}(\mathit{window}(M),\phi)$\;
  $E\gets$ the safe prefix ending at the first non-exhausted child window that
  is depleted\;
  Move $E$ from the window to $\mathit{ready}(M)$ and retain the uncertified
  suffix\;
}
$X\gets$ up to $c$ items from the front of $\mathit{ready}(M)$\;
$\mathit{delivered}(M)\gets\mathit{delivered}(M)+|X|$\;
\Return{$X$}\;
\end{algorithm}

\begin{algorithm}[t]
\small
\caption{$\textsc{Next}(M,c)$ for a Loser-Tree Producer}
\label{alg:pipelined_lt_next}
\KwIn{binary or segmented loser-tree producer $M$, request $c$,
      instruction $\phi$, batch capacity $m$, segment length $h$ when
      segmented, budget $K$}
\KwOut{up to $c$ next items}

$c\gets\min(c,K-\mathit{delivered}(M))$; $X\gets[\ ]$\;
\If{$M$ is uninitialized}{
  Pull the initial head or segment from every child concurrently and
  then initialize the corresponding binary or segmented loser tree\;
}
\While{$|X|<c$ \textbf{and} $M$ has remaining output}{
  \uIf{$M$ is binary}{
    $u\gets$ the root winner; append its representative to $X$\;
    Pull one replacement from child $u$, then call \textsc{RefreshPath} on its
    leaf-to-root path\;
  }
  \Else{
    $r\gets\textsc{Root}(M)$\;
    \If{$\mathit{ready}(r)=\emptyset$}{
      $ok\gets\textsc{PointerGuidedPathRefresh}(M,\phi,m,h)$\;
      \lIf{\textbf{not} $ok$}{\textbf{break}}
    }
    Append $\textsc{TakeReady}(r,\min(h,c-|X|))$ to $X$\;
  }
}
$\mathit{delivered}(M)\gets\mathit{delivered}(M)+|X|$\;
\Return{$X$}\;
\end{algorithm}

```latex
\section{To be revised:Global Orderedness of Segmented Loser Tree}
\label{sec:appendix-segment-correctness}

For an ordered sequence $S$, let $\operatorname{Pref}_q(S)$ denote its
first $\min\{q,|S|\}$ items. We distinguish two properties: every emitted
segment must be the next prefix of the global order, and every input item
must be emitted exactly once.

\paragraph{Why path selection matters.}
A locally ordered replay does not by itself guarantee global order.
Consider $h=2$, $L=[1,4]$, and $R=[2,6]$. Their replay produces
\[
tmp=[1,2],\qquad \mathit{Loser}(v)=[4,6],
\]
and the left slot is exhausted first. Suppose its next segment begins
with $5$, whereas the next right segment is $[7,8]$. If the algorithm
incorrectly refreshes the right path, replaying $[4,6]$ with $[7,8]$
would emit $[4,6]$, placing $6$ before the unseen item $5$. The replay
batch is locally sorted, but it omits a candidate required for global
order. Hence correctness depends on proving that
$\mathit{Pointer}(v)$ selects exactly the child whose next segment may
precede part of $\mathit{Loser}(v)$.

A replay segment is associated with a \emph{logical child slot}.
This association identifies the child whose continuation is required
when the segment is exhausted; it does not require every item in the
segment to originate from that child. After a replay, each internal node
maintains the following invariant:

\begin{itemize}
  \item $\mathit{Loser}(v)$ is the ordered segment of exposed items not propagated to the parent;
  \item if $\mathit{Pointer}(v)=p$, the next segment from $p$ is the only unexposed segment that may precede an item in $\mathit{Loser}(v)$; and
  \item if $\mathit{Pointer}(v)=\textsc{null}$, the entire loser segment is certified and may be propagated without another ranking call.
\end{itemize}

\begin{lemma}[First-exhausted replay dependency]
\label{lem:segment-frontier-barrier}
Let $L$ and $R$ be valid $h$-item replay segments associated with the two
logical child slots of $v$, and let
\[
  Z=\operatorname{Sort}_{\rho}(L\uplus R)
\]
under a rank map consistent with $\prec_\phi$. Let $p$ be the slot whose
$h$-th item is encountered first when scanning $Z$, and let $o$ be the
other slot. Then
\[
  tmp=Z[1..h],
  \qquad
  \mathit{Loser}(v)=Z[h+1..2h],
  \qquad
  \mathit{Pointer}(v)=p
\]
satisfy the node invariant. In particular, $tmp$ is the next certified
$h$-item prefix, and refreshing any slot other than $p$ is insufficient
to certify the next winner in general.
\end{lemma}

\begin{proof}
Let $b$ be the position at which $p$ is first exhausted. Since both
inputs contain $h$ items,
\[
  h\le b\le 2h-1.
\]
Consequently, every item of $Z[1..h]$ occurs no later than the last
exposed item of $p$. Any subsequent item from $p$ must follow that last
item, and therefore cannot precede $Z[1..h]$. The other slot $o$ is not
exhausted at position $b$; after $b$, all remaining items of $Z$ come
from $o$. Its unseen continuation therefore follows the final item of
$Z$. Thus $Z[1..h]$ is certified against both continuations.

In contrast, the next item from $p$ may lie between $Z[b]$ and
$Z[b+1]$. It may therefore precede part of
$Z[h+1..2h]$, whereas an unseen item from $o$ cannot. Hence $p$ is the
unique unresolved dependency and must become $\mathit{Pointer}(v)$.

The loser segment may contain items originating from both subtrees, but
its suffix following position $b$ consists only of items from $o$ and
ends with the exposed frontier of $o$. It can therefore occupy the
logical slot associated with $o$ during the next replay. If the new
segment from $p$ is exhausted first, $p$ remains the unresolved
dependency. If the stored loser is exhausted first, the exposed frontier
of $o$ has been consumed and the dependency switches to $o$. This is
exactly the pointer update performed by
Algorithm~\ref{alg:replay_segment_node}.
\end{proof}

Algorithm~\ref{alg:replay_segment_node} presents the steady-state case
of two full segments. For a final short segment, its actual length is
used as the exhaustion threshold. A short segment has no continuation;
therefore, if it is exhausted first, the pointer moves to the other slot
when that slot has more output, and becomes \textsc{null} otherwise.

\begin{theorem}[Conditional global orderedness]
\label{thm:segment-global-order}
Fix a strict total order $\prec_\phi$ with deterministic tie-breaking.
Suppose that:
\begin{itemize}
  \item[(A1)] every input run is finite and sorted by $\prec_\phi$;
  \item[(A2)] every rank map produced during initialization or path
        refresh is the restriction of the same order $\prec_\phi$;
  \item[(A3)] every unexhausted producer eventually returns its next
        segment, with only its final segment allowed to contain fewer
        than $h$ items; and
  \item[(A4)] the tree depth is at most $d$ and $(d+1)h\le m$.
\end{itemize}
Then Algorithm~\ref{alg:segmented_loser_tree_merge} emits exactly the
sorted merge of its input runs. Any consumer that stops after $K$ items
receives the first $K$ items of the same global order.
\end{theorem}

\begin{proof}
Let
\[
  \mathcal S=
  \operatorname{Sort}_{\prec_\phi}
  (R_1\uplus\cdots\uplus R_k)
\]
be the desired global merge.

\paragraph{Initialization.}
At line~\ref{line:slt-init}, each leaf returns the next prefix of its
sorted run. By (A1), this segment is a certified prefix of that leaf.
The shared rank map $\rho_0$ orders all initial leaf segments according
to $\prec_\phi$ by (A2).

We apply structural induction to \textsc{BuildSubtree}. A leaf returns
its certified segment directly. At an internal node, the recursive calls
return the next winner segments of its left and right subtrees. By
Lemma~\ref{lem:segment-frontier-barrier},
\textsc{ReplaySegmentNode} returns their next merged prefix and stores a
loser/pointer pair satisfying the node invariant. Hence the value
returned by every recursive call is the next certified segment of its
subtree. In particular, the initial root $tmp$ is the first segment of
$\mathcal S$. The returned values are transient: only
$\mathit{Loser}(v)$ and $\mathit{Pointer}(v)$ remain at an internal node.

\paragraph{Pointer-guided refresh.}
After a root winner is emitted, let
\[
  \mathcal P=(u_0,u_1,\ldots,u_t),
  \qquad u_0=\textsc{Root}(T),
\]
be the path followed by Algorithm~\ref{alg:repair_segment_path}.
At every internal node $u_i$, the node invariant states that the child
selected by $\mathit{Pointer}(u_i)$ is the only child whose continuation
may precede part of $\mathit{Loser}(u_i)$. Following any other edge can
therefore omit a necessary candidate, as illustrated by the preceding
counterexample. Following the stored pointers reaches exactly the
unresolved frontier.

If $u_t$ is a leaf, \textsc{FillLeaf} supplies its next certified
segment. If it is an internal node with a null pointer, its loser is
already fully certified and can be promoted directly. In either case,
the terminal node supplies a valid initial $tmp$.

The refresh batch is
\[
  B=\{tmp\}\uplus
    \biguplus_{v\in\mathcal P}\{\mathit{Loser}(v)\}.
\]
By (A2), $\rho_B$ gives the order induced by $\prec_\phi$ on every item
participating in the refresh. Replay then proceeds bottom-up. Assume the
current $tmp$ is the next certified segment returned by the pointer
child of a path node $v$. The stored loser occupies the other logical
slot. Lemma~\ref{lem:segment-frontier-barrier} shows that replay returns
the next certified segment of $v$ and establishes the new loser/pointer
invariant at $v$. Induction from $u_t$ to $u_0$ therefore proves that the
returned root $tmp$ is the next globally certified segment. Nodes
outside $\mathcal P$ require no update because their continuations cannot
precede the stored loser at any path node.

\paragraph{Root output.}
Let $\Pi_s$ be the output after $s$ root segments and let $U_s$ be the
items not yet emitted. We maintain
\[
  \Pi_s=\operatorname{Pref}_{|\Pi_s|}(\mathcal S),
  \qquad
  tmp_s=
  \operatorname{Pref}_{|tmp_s|}
  \left(\operatorname{Sort}_{\prec_\phi}(U_s)\right).
  \tag{I}
\]
Initialization establishes (I) for $s=0$. If (I) holds, appending
$tmp_s$ gives
\[
  \Pi_{s+1}
  =\Pi_s\mathbin{\|}tmp_s
  =\operatorname{Pref}_{|\Pi_s|+|tmp_s|}(\mathcal S).
\]
The pointer-guided refresh proved above returns the next prefix of the
new remainder $U_{s+1}$, reestablishing (I). Thus every emitted segment
extends the same global sorted prefix.

\paragraph{Conservation, capacity, and termination.}
A child transfers ownership of every item in its returned $tmp$ to its
parent. Ranking changes only their order. At a replay node,
\[
  tmp\mathbin{\|}\mathit{Loser}(v)
  =\operatorname{Sort}_{\rho}(L\uplus R),
\]
so every replayed occurrence is assigned to exactly one transient winner
or one persistent loser. Promotion and output transfer ownership without
copying. Therefore no item is lost or duplicated.

A refresh batch contains at most one $h$-item terminal segment and one
loser of at most $h$ items from each of at most $d$ path levels. Hence
\begin{equation}
  |B|\le (d+1)h\le m.
  \label{eq:appendix-repair-capacity}
\end{equation}
Every nonempty root $tmp$ emits at least one item. By (A3), an
unexhausted terminal eventually supplies its next segment; when a final
short segment is consumed, the boundary rule removes its exhausted
dependency. Because all runs are finite, the process terminates only
after every source, transient winner, and stored loser is empty.
Conservation then gives $|\Pi|=|\mathcal S|$, while invariant (I) gives
$\Pi=\mathcal S$.

A pipelined consumer that stops after $K$ items merely truncates this
same sequence and therefore receives
$\operatorname{Pref}_K(\mathcal S)$.
\end{proof}

The theorem depends on Assumption~(A2): enforcing a valid permutation
inside each individual ranking call does not by itself guarantee
consistency between different calls.
```

\section{Complete Parameter Reference}
\label{sec:parameter-reference}

\begin{table}[H]
  \caption{Complete parameters used in the algorithms and cost models.}
  \label{tab:params-full}
  \small
  \centering
  \begin{tabularx}{\linewidth}{c>{\raggedright\arraybackslash}X}
    \toprule
    \textbf{Variable} & \textbf{Definition} \\
    \midrule
    $D$ & Input item collection \\
    $n$ & Input size; $n=|D|$ \\
    $\phi$ & Natural-language ordering instruction \\
    $m$ & Listwise batch size; max.\ \#items in one LLM call \\
    $K$ & Number of items to return under \textsf{LIMIT}~$K$ \\
    $\Pi$, $\Pi_K$ & Output semantic order and its top-$K$ prefix \\
    $k$ & Merge fan-in: max.\ \#runs merged per group, chosen by the cost model \\
    $g$ & Effective fan-in of one group, tree, or producer node ($g \le k$); $g_t=\min(r_t,k)$ in pass $t$, $g_v$ at node $v$ \\
    $r$ & Number of initial sorted runs; $r=\lceil n/m \rceil$ \\
    $r_t$ & Number of surviving runs after $t$ merge passes; $r_t=\lceil r/k^t \rceil$ \\
    $s$, $G_t$ & Numbers of groups in the algorithm and in pass $t$; $G_t=\lceil r_t/k \rceil$ \\
    $P(k)$ & Number of merge passes; $P(k)=\lceil \log_k r \rceil$ \\
    $A$ & Set of active (non-exhausted) runs in a merge \\
    $h$ & Per-run window quota in the sliding-window merge, $\lfloor m/|A| \rfloor$; segment length in the segmented loser tree, $\lfloor m/(1{+}d) \rfloor$ \\
    $d$ & Loser-tree depth; $d=\max(1,\lceil \log_2 g \rceil)$ \\
    $\mu_g$ & Expected \#items emitted per listwise call at effective fan-in $g$ \\
    $C$ & Total \#listwise calls in the full-ranking cost models ($C_{\mathrm{sw}}$, $C_{\mathrm{seg}}$, $C_{\mathrm{bin}}$) \\
    \bottomrule
  \end{tabularx}
\end{table}

\section{Full Listwise Ranking Prompt Template}
\label{sec:appendix-listwise-prompt}

We use the following OrderByLLM listwise prompt with vLLM constrained decoding.  Line breaks are added only for typesetting.

\begin{verbatim}
[SYSTEM]
You are a helpful agent. Think step by step.
Output a JSON object.
[USER]
You are given a criterion and a list of {items}
with {item} id at the front.
Ranking rules:
1. Use only the Criterion and the provided item
   content.
2. Apply hard constraints first: entity, time,
   location, number, and required attributes.
3. Prefer direct matches to the Criterion over
   prerequisite, neighboring, or loosely related
   topics.
4. For domain-relevance queries, rank core topics
   in the target domain above supporting or
   adjacent fields.
5. Break close cases by directness, specificity,
   support, and completeness.
6. Produce a best-effort total order; do not leave
   broad ties between close items.
7. Ignore input order, text length, tone, and
   formatting.
Rank the {items} from WORST to BEST according to
the Criterion.
Criterion: {criterion}
{items}:
{id_prefix}:0
{item_0}
...
{id_prefix}:{num_items_minus_one}
{item_last}
Output only {"ranking":[...]} with IDs ordered
from WORST to BEST. Use every integer from 0 to
{num_items_minus_one} exactly once.
\end{verbatim}

For $m$ items, the vLLM processor follows fixed JSON tokens and an ID trie, uses a bit mask to prevent repeated IDs, and assigns $-\infty$ to invalid token logits.  It closes the object after all $m$ IDs appear, producing a permutation of $0,\ldots,m-1$.

\section{Alternative cost-surface layout (integration appendix)}
\label{sec:appendix-e2-2x3}

% Approved: PR2 — 2x3 appendix layout variant
\begin{figure*}[t]
  \centering
  \includegraphics[width=\textwidth]{e2-layout-2x3.pdf}
  \caption{Synthetic comparison layout, not experimental results.
  Three query slots by Full and LIMIT10, LT only, same camera and marker rules as the $2\times 2$ layout in the main text (\cref{fig:exp2-surface-2x2}).
  All six panels share the LT legal window and one $z$ scale.
  Query slots are labels.}
  \label{fig:exp2-surface-2x3}
\end{figure*}
